% !TeX root = ../papers/nc-p3.tex
\section{素体の二つの放物線に沿う全剰余近傍}
\label{sec:rank-fifteen-prime-parabolic-pencils}
\begin{theorem}\label{thm:rank-fifteen-prime-parabolic-pencils}
剰余標数3の付値で $X=\bar x\ne0,1$ とし、$Y=\bar y$ は非零とする。
$Y/(1-X)^2\in\{1,-1\}$ なら全ての通常な標数零の階数15窓が
標準Koszul左右AS正則で、各始点が四中心四次式上自由階数256である。
$X$ の剰余次数に制限はない。
\end{theorem}
\begin{proof}
$Y=(1-X)^2$ では $c^2=X$, $v=1/c$ を取る。$K=c$, $H=K^2-1$ となり、
六根のうち $p_0^2=(K^2-1+2H)/4=0$ である。他の五根は $c\ne0,\pm1$ により非零である。
第\ref{sec:rank-fifteen-genus-two-unit-line}節の36成分から $p_0$ の四成分を除き、
残る32語に同じ対角変換の規格化を施す。
\path{computations/verify_rank_fifteen_prime_parabolic_pencils.py} と共通の規格化コードは、
指数行列の独立な13行の小行列式が絶対値4で、全ての規格化係数で五根の指数が偶数であることを確かめる。
従って $\mathbb F_3(c)$ で四窓を扱える。変換後の四次中心は文字0に五次元あり、
後ろの四自由列を選んだ四中心商が四始点で十三次まで標準次元に一致する。
始点別の除外多項式の次数は $7,3,5,5$、最小公倍式は
\[
 c(c-1)(c+1)(c^2+1)(c^2+c+2)
\]
である。$c,-c$ の両方が零点となるのは最初の四因子に限られる。
仮定のもとで残る $X=-1$ は前に認証した。

一般に $R=Y/(1-X)^2$ を固定し、$R\ne0,1$ とする。
$r^2=R$, $a^3=r$, $b=1/c$, $d=b^3-b$ とし、$v^3-v=r d$ の符号を選ぶ。
$q=v-ab$ と置くと
\[
 b=(q-q^3)/(r-a),\qquad v=q+ab,\qquad Z=\sqrt{R-1}\,d .
\]
$r-a\ne0$ であるから、この一変数表示は全ての $X\ne0,1$ に使える。
$Z^2=(v^3-v)^2-d^2$ であり、36成分の正規化に現れる六根の積も $q$ の有理関数になる。

$R=-1$ では $r=i$, $a=-i$, $Z=d$ とできる。$1/X=b^2=-(q^3-q)^2$ なので、
各有限非零水準の六表示は $q\mapsto\pm q+j$ ($j\in\mathbb F_3$) で尽くされる。
$\mathbb F_9(q)$ で同じ中心と四中心商を計算し、逆元を素体へのノルムで記録すると
四始点の除外次数は $1174,1285,1375,1074$、最小公倍式の次数は4438である。
$1/X=0,1$ を除き六変換で安定な部分は次数126で、完全ファイバーとなる
$1/X$ の閉点の次数は $1,2,2,2,3,3,4,4$ となる。
二次までは前節で認証済みである。二つの三次と二つの四次では各々
$\mathbb F_{3^{12}},\mathbb F_{3^8}$ で元の倍率を取り、四始点の中心商を十三次まで検算する。
$Y=-(1-X)^2$ は素体上の式なので、各多項式につき一根でFrobenius軌道が尽くされる。
両水準で単元主元による持上げと全次数判定は第\ref{sec:rank-fifteen-genus-two-unit-line}節による。
\end{proof}

\begin{corollary}
標数零の代数閉体で $\lambda\in\mathbb Q^\times$ とする。
放物線 $y=\lambda(1-x)^2$ の非零水準の通常窓は全て上の正則性と自由性を持つ。
\end{corollary}
\begin{proof}
3進付値を係数体に延長する。負付値の $x$ と $\bar x=0,1$ は全水準で認証済みである。
残る $x$ では $1-x$ が単元である。$\lambda$ が非単元なら $y$ も非単元で、
既知の内側または極の定理を使う。単元なら $\bar\lambda=\pm1$ なので本定理を使う。
\end{proof}
